\subsection{终拓扑}

命题 6. 设 X 是一个集合，设 \((\mathbf{Y}_t)_{t \in I}\) 是拓扑空间的一个族，并且对每个 \(t \in I\)，设 \(f_t\) 是由 \(\mathbf{Y}_t\) 到 X 中的映射。设 \(\mathfrak{D}\) 是 X 的这样的子集 U 所成之集：\(\overline{f}_t^1(\mathbf{U})\) 在 \(\mathbf{Y}_t\) 中是开集（对每个 \(t \in I\)）；则 \(\mathfrak{D}\) 是 X 上某个拓扑 \(\mathfrak{C}\) 的开子集所成之集，该拓扑是 X 关于族 \((f_t)\) 的终结构（《集合论》，第四章，§ 2，第 5 目），特别地，\(\mathfrak{C}\) 是使诸映射 \(f_t\) 连续的 X 上最细的拓扑。换言之，若 g 是由 X 到拓扑空间 Z 中的映射，则 g 连续（X 赋以拓扑 \(\mathfrak{C}\)）当且仅当每个映射 \(g \circ f_t\) 连续。

立即可以验证 \(\mathfrak{D}\) 满足公理 \((\mathrm{O_I})\) 与 \((\mathrm{O_{II}})\)。{[}Set Theory, R, § 4, formulae (34) and (46){]}。我们将证明命题的最后一个断言，由终结构的一般性质（《集合论》，第四章，§ 2，第 5 目，判别准则 CST 18），它蕴含其余断言。显然诸 \(f_{i}\) 在拓扑 \(\mathfrak{C}\) 中连续，这由 \(\mathfrak{D}\) 的定义得出（第 1 目，定理 1）；因此若 \(g\) 连续，则每个映射 \(g \circ f_{i}\) 连续（第 1 目，定理 2）。反之，设每个 \(g \circ f_{i}\) 都连续，并设 V 是 Z 中的开集；由假设，\(\overline{f}_{i}^{1}(\overline{g}^{1}(V))\) 在 \(Y_{i}\) 中是开集（对每个 \(i \in I\)）；于是 \(\overline{g}^{1}(V) \in \mathfrak{D}\)，证明完成。

推论. 在命题 6 的假设下，X 的子集 F 在拓扑 \(\mathfrak{C}\) 中是闭集当且仅当 \(\overline{f}_{\iota}^{1}(\mathbf{F})\) 在 \(Y_{\iota}\) 中是闭集（对每个 \(\iota \in I\)）。这由拓扑 \(\mathfrak{C}\) 中开集的定义取余集得出。

终结构的一般性质（《集合论》，第四章，§ 2，第 5 目，判别准则 CST 19）蕴含下列传递性质（其直接证明同样是容易的）：

命题 7. 设 X 是一个集合，\((Z_{\iota})_{\iota \in I}\) 是拓扑空间的一个族，\((J_{\lambda})_{\lambda \in L}\) 是 I 的一个分拆，而 \((Y_{\lambda})_{\lambda \in L}\) 是以 L 为指标的一族集合。又对每个 \(\lambda \in L\)，设 \(h_{\lambda}\) 是由 \(Y_{\lambda}\) 到 X 中的一个映射；对每个 \(\lambda \in L\) 与每个 \(\iota \in J_{\lambda}\)，设 \(g_{\lambda \iota}\) 是由 \(Z_{\iota}\) 到 \(Y_{\lambda}\) 中的一个映射，并令 \(f_{\iota} = h_{\lambda} \circ g_{\lambda \iota}\)。如果每个 \(Y_{\lambda}\) 都赋以使诸映射 \(g_{\lambda \iota} (\iota \in J_{\lambda})\) 连续的最细拓扑，那么使诸 \(f_{\iota}\) 连续的 X 上最细的拓扑与使诸 \(h_{\lambda}\) 连续的最细拓扑相同。

\paragraph*{例}

\begin{enumerate}
\def\labelenumi{\arabic{enumi})}
\item
  商拓扑。设 X 是拓扑空间，R 是 X 上的一个等价关系，Y = X/R 是 X 关于关系 R 的商集，\(\varphi: X \to Y\) 是典范映射。使 \(\varphi\) 连续的 Y 上最细的拓扑称为 X 的拓扑被关系 R 的商；我们将在 § 3 中更详细地研究它。
\item
  拓扑集的下确界。拓扑的每个族 \((\mathfrak{C}_t)_{i \in I}\)（集合 \(X\) 上的拓扑的族）都有一个下确界 \(\mathfrak{C}\)（在 \(X\) 上所有拓扑的集合中），即 \(\mathfrak{C}\) 是 \(X\) 上比每个 \(\mathfrak{C}_t\) 都粗的所有拓扑中最细的一个。为看出这一点，可以应用命题 6，取 \(Y_t\) 为集合 \(X\)（赋以拓扑 \(\mathfrak{C}_t\)），取 \(f_t\) 为恒等映射 \(Y_t \to X\)。若 \(\mathfrak{D}_t\) 是 \(X\) 的在拓扑 \(\mathfrak{C}_t\) 中为开集的子集所成之集，则集合 \(\bigcap_{i \in I} \mathfrak{D}_t\) 是 \(X\) 的在 \(\mathfrak{C}\) 中为开集的子集所成之集。
\end{enumerate}

\(\mathfrak{G}\) 也称为诸拓扑 \(\mathfrak{G}_{t}\) 的交。

\begin{enumerate}
\def\labelenumi{\arabic{enumi})}
\setcounter{enumi}{2}
\tightlist
\item
  拓扑空间的和。设 \((\mathbf{X}_t)_{i \in I}\) 是拓扑空间的一个族，\(\mathbf{X}\) 是诸 \(\mathbf{X}_t\) 的和所成的集合（《集合论》，第二章，§ 4，第 8 目，定义 8）；对每个 \(i \in I\)，设 \(j_t\) 是由 \(\mathbf{X}_t\) 到 \(\mathbf{X}\) 中的典范（单射）映射。最细拓扑 \(\mathfrak{G}\)（\(\mathbf{X}\) 上的、使诸映射 \(j_t\) 连续者）称为诸 \(\mathbf{X}_t\) 的拓扑之和，而 \(\mathbf{X}\) 赋以这个拓扑后称为诸拓扑空间 \(\mathbf{X}_t\) 的和。让我们把每个 \(\mathbf{X}_t\) 与 \(\mathbf{X}\) 的一个子集借助 \(j_t\) 等同起来；则集合 \(\mathbf{A} \subset \mathbf{X}\) 在拓扑 \(\mathfrak{G}\) 中是开（相应地，闭）集当且仅当诸集合 \(\mathbf{A} \cap \mathbf{X}_t\) 中的每一个在 \(\mathbf{X}_t\) 中都是开（相应地，闭）集。特别地，诸 \(\mathbf{X}_t\) 中的每一个都既开又闭。
\end{enumerate}

下列命题推广了例 3 的情形：

命题 8. 设 X 是一个集合，\((\mathbf{X}_{\lambda})_{\lambda \in \mathbf{L}}\) 是 X 的子集的一个族。假设每个 \(X_{\lambda}\) 都赋有一个拓扑 \(C_{\lambda}\)，使得对每对指标 \((\lambda, \mu)\)：

\begin{enumerate}
\def\labelenumi{\arabic{enumi})}
\item
  \(\mathbf{X}_{\lambda} \cap \mathbf{X}_{\mu}\) 在拓扑 \(\mathfrak{C}_{\lambda}, \mathfrak{C}_{\mu}\) 的每一个中都是开（相应地，闭）集。
\item
  在 \(\mathbf{X}_{\lambda} \cap \mathbf{X}_{\mu}\) 上由 \(\mathfrak{C}_{\lambda}\) 与 \(\mathfrak{C}_{\mu}\) 诱导的拓扑相重合。设 \(\mathfrak{C}\) 是 \(\mathbf{X}\) 上使诸单射 \(j_{\lambda}: \mathbf{X}_{\lambda} \to \mathbf{X}\) 连续的最细拓扑。则对每个 \(\lambda \in \mathbf{L}\)，\(\mathbf{X}_{\lambda}\) 在 \(\mathbf{X}\) 中在拓扑 \(\mathfrak{C}\) 下是开（相应地，闭）集，并且 \(\mathfrak{C}\) 在 \(\mathbf{X}_{\lambda}\) 上诱导的拓扑与 \(\mathfrak{C}_{\lambda}\) 相重合。
\end{enumerate}

鉴于命题 6 及其推论，只需证明：对每个 \(\lambda\) 与每个子集 \(A_{\lambda}\)（\(X_{\lambda}\) 的子集），下列命题等价：

\begin{enumerate}
\def\labelenumi{(\roman{enumi})}
\item
  \(\mathbf{A}_{\lambda}\) 在拓扑 \(\mathfrak{C}_{\lambda}\) 中是开（相应地，闭）集。
\item
  对每个 \(\mu \in L\)，\(A_{\lambda} \cap X_{\mu}\) 在拓扑 \(\mathfrak{G}_{\mu}\) 中是开（相应地，闭）集。
\end{enumerate}

显然，取 \(\mu = \lambda\) 即可看出 (ii) 蕴含 (i)。反之，若 (i) 成立，则 \(A_{\lambda} \cap X_{\mu}\) 在 \(X_{\lambda} \cap X_{\mu}\) 中相对于拓扑 \(\mathcal{C}_{\lambda\mu}\) 是开（相应地，闭）集；该拓扑在 \(X_{\lambda} \cap X_{\mu}\) 上由 \(\mathcal{C}_{\lambda}\) 诱导；但 \(\mathcal{C}_{\lambda\mu}\) 也是在 \(X_{\lambda} \cap X_{\mu}\) 上由 \(\mathcal{C}_{\mu}\) 诱导的拓扑；因此 \(A_{\lambda} \cap X_{\mu}\) 也是 \(X_{\lambda} \cap X_{\mu}\) 与一个子集 \(B_{\mu}\) 的交，其中该子集是 \(X_{\mu}\) 的子集且在拓扑 \(\mathcal{C}_{\mu}\) 中是开（相应地，闭）集；由于 \(X_{\lambda} \cap X_{\mu}\) 在 \(\mathcal{C}_{\mu}\) 中是开（相应地，闭）集，\(A_{\lambda} \cap X_{\mu}\) 也如此。证明完毕。

注意，若诸 \(\mathbf{X}_{\lambda}\) 的并不同于 \(\mathbf{X}\)，则 \(\mathcal{G}\) 在 \(\mathbf{X} - \left(\bigcup_{\lambda \in \mathbf{L}} \mathbf{X}_{\lambda}\right)\) 上诱导的拓扑是离散的。这是因为，若 \(x \in \mathbf{X}\) 不属于诸 \(\mathbf{X}_{\lambda}\) 中的任何一个，则 \(\{x\} \cap \mathbf{X}_{\lambda} = \emptyset\) 在每个拓扑 \(\mathcal{G}_{\lambda}\) 中都是开集，因而 \(\{x\}\) 在拓扑 \(\mathcal{G}\) 中是开集。

